\section{A shorter infrared coefficient}
\label{sec:comp-ir-simplification}
The two roots in \eqref{eq:comp-soft-integrated-ir} are reciprocal.
This relation reduces its endpoint sum to one dilogarithmic function.
Set
\begin{equation}
 X=s-1>0,\qquad Y=\sqrt{\frac{X}{X+4}},\qquad
 q=\frac{X+2-\sqrt{X(X+4)}}2,\qquad L=\ln(1+X).
\end{equation}
Then $0<q<1$, $X=(1-q)^2/q$, and $z_-=q$, $z_+=q^{-1}$.
Define
\begin{equation}
 \mathcal G(X)=\frac12\ln^2q
 +\frac13\operatorname{Li}_2(-q^3)
 -\operatorname{Li}_2(-q)-\frac{\pi^2}{18}.
 \label{eq:comp-ir-short-function}
\end{equation}
Both dilogarithm arguments lie between $-1$ and zero.  All functions
in this expression are consequently real on the physical interval,
with their branches fixed without a continuation across a cut.
The integrated infrared coefficient is
\begin{equation}
 \begin{aligned}
 A(1+X)={}&-\frac{2(X^2-4X-6)}{X^3}L^2
 -\frac{2P(X)}{X^3(1+X)^2(X+4)}L\\
 &+\frac{12YQ(X)}{X^4(X+4)}\mathcal G(X)
 -\frac{R(X)}{X^2(1+X)^2},
 \end{aligned}
 \label{eq:comp-ir-short}
\end{equation}
where
\begin{align}
 P(X)&=X^5-22X^4-184X^3-444X^2-424X-144,\\
 Q(X)&=X^4+2X^3-24X^2-80X-48,\\
 R(X)&=5X^3+60X^2+100X+48.
\end{align}
The normalization is $A(s)=\mathcal C_{\mathrm{IR}}(s)/s$, so
$A(s)\ln(\lambda_{\mathrm{phys}}/m)$ is the logarithmic term of the
complete virtual coefficient in units of $\alpha_{\mathrm T}^3/m^2$.
Equation~\eqref{eq:comp-ir-short} is exact for every $X>0$.
Its apparent threshold poles cancel in the combined expression.

To derive it directly from the angular integral, denote the endpoint
difference of a primitive by $\Delta J=J(s)-J(1)$.  The two logarithmic
pole primitives in \eqref{eq:comp-soft-integrated-ir} satisfy
\begin{align}
 \Delta H_1(q)&=-\ln q, & \Delta H_1(q^{-1})&=\ln q,\\
 \Delta H_2(q)&=1, & \Delta H_2(q^{-1})&=1.
\end{align}
Their residues obey $c_{-,q}=c_{-,q^{-1}}$ and
$d_{-,q}+d_{-,q^{-1}}=8$.  Thus the logarithms of the negative-sign
sector cancel, while its double poles give the constant eight.
For the positive-sign sector, $c_{+,q}=-c_{+,q^{-1}}$.
Use the real dilogarithm inversion identity
\begin{equation}
 \operatorname{Re}\operatorname{Li}_2(u)
 +\operatorname{Re}\operatorname{Li}_2(u^{-1})
 =\frac{\pi^2}{3}-\frac12\ln^2u,\qquad u>0,
\end{equation}
to combine its simple-pole endpoint functions into
\begin{equation}
 \begin{aligned}
 \Delta J_1(q)-\Delta J_1(q^{-1})
 ={}&2\operatorname{Li}_2(q)-\operatorname{Li}_2(q/s)
  -\operatorname{Li}_2(qs)\\
 &-2L\ln q-\frac12L^2\equiv\mathcal F(q).
 \end{aligned}
\end{equation}
Here $qs=1-q+q^2<1$ and $q/s=q^2/(1-q+q^2)<1$.
Differentiation and the identity
$(1+q)(1-q+q^2)=1+q^3$ give
\begin{equation}
 \frac{\dd\mathcal F}{\dd q}=-\frac{3L}{q},\qquad
 \frac{\dd\mathcal G}{\dd q}=-\frac{L}{q},\qquad
 \mathcal F(1)=\mathcal G(0)=0.
\end{equation}
In the last boundary value the argument of $\mathcal G$ is $X$,
which tends to zero as $q\to1$.  Therefore
$\mathcal F(q)=3\mathcal G(X(q))$, with no undetermined constant.
The double-pole primitives supply only elementary terms,
\begin{align}
 \Delta J_2(q)&=-\frac{L}{s-q}-\frac{L+\ln q}{q},\\
 \Delta J_2(q^{-1})&=-\frac{L}{s-q^{-1}}+q(\ln q-L).
\end{align}
Substituting these differences and the elementary power primitives
into \eqref{eq:comp-soft-integrated-ir} yields
\eqref{eq:comp-ir-short}.  This derivation uses only the original
Born-weighted soft angular integral.

An independent internal check differentiates the complete virtual
23-word result with respect to $\ln(\lambda_{\mathrm{phys}}/m)$.
Its four nonzero coefficients agree algebraically with
\eqref{eq:comp-ir-short}, using
$G(-1;X)=L$, $G(-1,-1;X)=L^2/2$ and
$G(r,-1;X)=\mathcal G(X)$.  This check does not set any integration
constant in the soft calculation.
